\documentclass[tikz,border=0mm]{standalone}
\usepackage{xcolor}
\usetikzlibrary{calc}

\begin{document}

\definecolor{minimizergreen}{RGB}{0,150,70}
\definecolor{overlapred}{RGB}{210,20,20}

\newcommand{\minimizer}[1]{\textcolor{minimizergreen}{#1}}
\newcommand{\nucStep}{0.44}

\tikzset{
  nuc n/.style={
    font=\ttfamily\large,
    inner sep=1pt,
    outer sep=0pt
  },
  nuc m/.style={
    font=\ttfamily\large,
    inner sep=1pt,
    outer sep=0pt,
    text=minimizergreen
  },
  idx/.style={
    font=\scriptsize,
    text=black!55,
    inner sep=1pt
  },
  strand label/.style={
    font=\small\itshape,
    text=black!70,
    anchor=east
  },
  section title/.style={
    font=\normalsize\bfseries,
    align=center,
    text width=3.5cm
  },
  section note/.style={
    font=\small,
    align=center,
    text width=3.5cm
  },
  bad note/.style={
    font=\small,
    align=center,
    text width=3.5cm,
    text=overlapred
  },
  good note/.style={
    font=\small,
    align=center,
    text width=3.5cm,
    text=minimizergreen
  },
  frame bad/.style={
    draw=overlapred,
    rounded corners=2pt,
    line width=0.7pt
  },
  frame good/.style={
    draw=minimizergreen,
    rounded corners=2pt,
    line width=0.7pt
  },
  col bad/.style={font=\small, text=overlapred, anchor=north},
  col good/.style={font=\small, text=minimizergreen, anchor=north}
}

% ------------------------------------------------------------
% One nucleotide word
%   #1 node id prefix, #2 x shift, #3 y shift,
%   #4 left-hand label, #5 list of letter/kind (n = plain, m = minimizer)
% ------------------------------------------------------------
\newcommand{\drawWord}[5]{%
  \begin{scope}[shift={(#2,#3)}]
    \node[strand label] at (-0.28,0) {#4};
    \foreach \b/\kind [count=\i from 0] in {#5}{%
      \node[nuc \kind] (#1\i) at ({\i*\nucStep},0) {\b};
    }%
  \end{scope}%
}

% Position ruler under a word: #1 x, #2 y, #3 last index
\newcommand{\drawRuler}[3]{%
  \begin{scope}[shift={(#1,#2)}]
    \foreach \i in {0,...,#3}{%
      \node[idx] at ({\i*\nucStep},0) {\i};
    }%
  \end{scope}%
}

% Box the occurrence spanning letters #2..#3 of word #1 and label it #4
%   #5 = style suffix (bad / good)
\newcommand{\markOcc}[5]{%
  \draw[frame #5]
    ($(#1#2.north west)+(-0.05,0.06)$) rectangle ($(#1#3.south east)+(0.05,-0.06)$);
  \node[col #5] at ($(#1#2.south east)+(0,-0.16)$) {#4};
}

\begin{tikzpicture}[font=\sffamily]

% ============================================================
% 1. the same k-mer seen from both strands
% ============================================================
\node[section title] at (0.95,1.75) {1. One $k$-mer, two strands};

\drawWord{Af}{0}{0.75}{forward}{A/m, T/m, C/n, G/n, A/m, T/m, G/n}
\drawRuler{0}{0.36}{6}
\drawWord{Ar}{0}{-0.75}{rev.\ comp.}{C/n, A/m, T/m, C/n, G/n, A/m, T/m}
\drawRuler{0}{-1.14}{6}

\node[section note, anchor=north] at (0.95,-1.65)
  {$k=7$, $m=2$: the minimal $2$-mer \minimizer{\texttt{AT}} is its own reverse
   complement and occurs twice, at mirrored positions};

% ============================================================
% 2. leftmost occurrence -- two different columns
% ============================================================
\node[section title] at (5.02,1.75) {2. Leftmost occurrence};

\drawWord{Bf}{3.7}{0.75}{}{A/m, T/m, C/n, G/n, A/m, T/m, G/n}
\markOcc{Bf}{0}{1}{column $5$}{bad}

\drawWord{Br}{3.7}{-0.75}{}{C/n, A/m, T/m, C/n, G/n, A/m, T/m}
\markOcc{Br}{1}{2}{column $1$}{bad}

\node[bad note, anchor=north] at (5.02,-1.65)
  {each strand takes the leftmost occurrence \emph{it} sees, so one $k$-mer
   reaches two columns};

% ============================================================
% 3. most central occurrence -- one column
% ============================================================
\node[section title] at (8.72,1.75) {3. Most central occurrence};

\drawWord{Cf}{7.4}{0.75}{}{A/m, T/m, C/n, G/n, A/m, T/m, G/n}
\markOcc{Cf}{4}{5}{column $1$}{good}

\drawWord{Cr}{7.4}{-0.75}{}{C/n, A/m, T/m, C/n, G/n, A/m, T/m}
\markOcc{Cr}{1}{2}{column $1$}{good}

\node[good note, anchor=north] at (8.72,-1.65)
  {centrality is mirror-symmetric, so both strands frame the same occurrence};

\end{tikzpicture}
\end{document}
